\section{Conclusion}
\label{sec:conclusion}

Do world models make better robots? On the evidence of the benchmarks the field has built, the honest
answer is that we cannot yet tell. The question asks for a closed-loop, per-capability comparison
between a world-model policy and a direct VLA policy, and our catalogue of 160 web-verified benchmarks
shows that almost none is built to run it: 138 of 160 are model-agnostic, only 11 build an explicit
contrast, counterfactual capability is nearly unmeasured, and only four benchmarks carry prediction
into executed action. The obstacle is methodological. World-model suites score prediction without
executing it, task-success suites execute a single policy without contrasting families, and the term
``world model'' itself names three different objects, none of which is action-conditioned,
closed-loop, and VLA-comparable at once.

This survey's contribution is to make that gap precise and measurable. We give an operational
taxonomy of the landscape across four evaluation lanes, a coverage comparison showing that no prior
survey crosses capability with model family, an evaluation loop that isolates the advantage of
prediction, and a four-metric protocol, anchored on named testbeds, that any benchmark could adopt to
report where and by how much prediction helps. The next benchmark that matters will not be the one
with the most tasks or the most realistic renders; it will be the one that runs both branches in one
loop and reads off the difference. Building it is the empty cell this survey names.
